% Options for packages loaded elsewhere
\PassOptionsToPackage{unicode}{hyperref}
\PassOptionsToPackage{hyphens}{url}
%
\documentclass[
  ignorenonframetext,
]{beamer}
\usepackage{pgfpages}
\setbeamertemplate{caption}[numbered]
\setbeamertemplate{caption label separator}{: }
\setbeamercolor{caption name}{fg=normal text.fg}
\beamertemplatenavigationsymbolsempty
% Prevent slide breaks in the middle of a paragraph
\widowpenalties 1 10000
\raggedbottom
\setbeamertemplate{part page}{
  \centering
  \begin{beamercolorbox}[sep=16pt,center]{part title}
    \usebeamerfont{part title}\insertpart\par
  \end{beamercolorbox}
}
\setbeamertemplate{section page}{
  \centering
  \begin{beamercolorbox}[sep=12pt,center]{part title}
    \usebeamerfont{section title}\insertsection\par
  \end{beamercolorbox}
}
\setbeamertemplate{subsection page}{
  \centering
  \begin{beamercolorbox}[sep=8pt,center]{part title}
    \usebeamerfont{subsection title}\insertsubsection\par
  \end{beamercolorbox}
}
\AtBeginPart{
  \frame{\partpage}
}
\AtBeginSection{
  \ifbibliography
  \else
    \frame{\sectionpage}
  \fi
}
\AtBeginSubsection{
  \frame{\subsectionpage}
}
\usepackage{amsmath,amssymb}
\usepackage{iftex}
\ifPDFTeX
  \usepackage[T1]{fontenc}
  \usepackage[utf8]{inputenc}
  \usepackage{textcomp} % provide euro and other symbols
\else % if luatex or xetex
  \usepackage{unicode-math} % this also loads fontspec
  \defaultfontfeatures{Scale=MatchLowercase}
  \defaultfontfeatures[\rmfamily]{Ligatures=TeX,Scale=1}
\fi
\usepackage{lmodern}
\ifPDFTeX\else
  % xetex/luatex font selection
\fi
% Use upquote if available, for straight quotes in verbatim environments
\IfFileExists{upquote.sty}{\usepackage{upquote}}{}
\IfFileExists{microtype.sty}{% use microtype if available
  \usepackage[]{microtype}
  \UseMicrotypeSet[protrusion]{basicmath} % disable protrusion for tt fonts
}{}
\makeatletter
\@ifundefined{KOMAClassName}{% if non-KOMA class
  \IfFileExists{parskip.sty}{%
    \usepackage{parskip}
  }{% else
    \setlength{\parindent}{0pt}
    \setlength{\parskip}{6pt plus 2pt minus 1pt}}
}{% if KOMA class
  \KOMAoptions{parskip=half}}
\makeatother
\usepackage{xcolor}
\newif\ifbibliography
\setlength{\emergencystretch}{3em} % prevent overfull lines
\providecommand{\tightlist}{%
  \setlength{\itemsep}{0pt}\setlength{\parskip}{0pt}}
\setcounter{secnumdepth}{-\maxdimen} % remove section numbering
\usepackage{arev} % order matters, putting this above allows FiraSans to override it for body text
\usepackage[sfdefault]{FiraSans}
\usepackage{inconsolata}
\usepackage{graphicx}
\usepackage{hyperref}
\usepackage{listings}
\usepackage{tikz}
\usepackage{pgfplots}
\usepackage{amsmath}
\usepackage[most]{tcolorbox}

\usepackage{xcolor}
\definecolor{blackish}{RGB}{56,58,54}
\definecolor{redish}{RGB}{109,41,49}
\definecolor{red}{RGB}{152,41,50}
\definecolor{orangeish}{RGB}{188,71,0}
\definecolor{blueish}{RGB}{25,33,139}

\usetheme{default}
\setbeamertemplate{navigation symbols}{}
\setbeamertemplate{itemize items}[triangle]
\setbeamertemplate{enumerate items}[default]

\setbeamercolor{structure}{fg=blueish}
\setbeamercolor{title}{fg=redish}
\setbeamercolor{subtitle}{fg=blackish}
\setbeamercolor{frametitle}{fg=blackish}
\setbeamercolor{frametitle rule}{fg=redish}
\setbeamercolor{footline}{fg=blackish}
\setbeamercolor{author}{fg=blackish}
\setbeamercolor{date}{fg=blackish}
\setbeamercolor{normal text}{fg=blackish}
\setbeamercolor{block title}{fg=white,bg=blueish!80}
\setbeamercolor{block body}{bg=blueish!5}
\setbeamercolor{block title alerted}{fg=white,bg=red!80}
\setbeamercolor{block body alerted}{bg=red!5}
\setbeamercolor{block title example}{fg=white,bg=orangeish!80}
\setbeamercolor{block body example}{bg=orangeish!5}

\setbeamertemplate{frametitle}{
  \vskip0.2em
  {\usebeamerfont{frametitle}\usebeamercolor[fg]{frametitle}\insertframetitle\par}
  {\color{redish}\hrule height 0.5pt \vspace{0.1em}}
  \vskip0.2em
}

\setbeamertemplate{footline}{%
  \hfill%
  {\color{red}\scriptsize COMP3000 Programming Languages}\hspace{0.4em}%
  {
\includegraphics[height=0.4cm]{logo.jpg}}\hspace{0.4em}%
  {\color{blackish}\scriptsize\insertframenumber/\inserttotalframenumber}\hspace{0.4em}\vspace{0.2em}
}

\setbeamertemplate{title page}{%
  \vbox{}
  \vfill
  \begingroup
    \centering
    {\usebeamerfont{title}\usebeamercolor[fg]{title}\inserttitle\par}
    \vskip0.5em
    {\usebeamerfont{subtitle}\usebeamercolor[fg]{subtitle}\insertsubtitle\par}
    \vskip1em
    {\color{redish}\hrule height 0.8pt\par}
    \vskip1em
    {\usebeamerfont{author}\usebeamercolor[fg]{author}\insertauthor\par}
    \vskip0.3em
    {\usebeamerfont{date}\usebeamercolor[fg]{date}\insertdate\par}
  \endgroup
  \vfill
}

\DeclareOldFontCommand{\bf}{\normalfont\bfseries}{\mathbf}
\DeclareOldFontCommand{\tt}{\normalfont\ttfamily}{\texttt}

\lstset{basicstyle=\ttfamily}

\newtcolorbox{aside}[1][]{
  title=Aside,
  width=0.9\textwidth,
  fonttitle=\bfseries,
  breakable,
  fonttitle=\bfseries\color{black},
  colframe=blueish!80,
  colback=blueish!2
  #1}

\newtcolorbox{goodtoknow}[1][]{
  title=Good to know,
  width=0.9\textwidth,
  fonttitle=\bfseries,
  breakable,
  fonttitle=\bfseries\color{black},
  colframe=blueish!80,
  colback=blueish!2
  #1}

\newtcolorbox{notebox}[1][]{
  title=Note,
  width=\textwidth,
  fonttitle=\bfseries,
  breakable,
  fonttitle=\bfseries\color{black},
  colframe=orangeish!80,
  colback=orangeish!2
  #1}

\newtcolorbox{hint}[1][]{
  title=Hint,
  width=\textwidth,
  fonttitle=\bfseries,
  breakable,
  fonttitle=\bfseries\color{white},
  colframe=blueish!80,
  colback=blueish!2
  #1}

\newtcolorbox{todo}[1][]{
  title=!! TODO !!,
  width=\textwidth,
  fonttitle=\bfseries,
  breakable,
  fonttitle=\bfseries\color{white},
  colframe=red!80,
  colback=red!2
  #1}

\providecommand{\tightlist}{%
  \setlength{\itemsep}{0pt}\setlength{\parskip}{0pt}}

\hypersetup{
  colorlinks=true,
  linkcolor=blueish,
  urlcolor=blueish,
  filecolor=blueish
}

\ifLuaTeX
  \usepackage{selnolig}  % disable illegal ligatures
\fi
\usepackage{bookmark}
\IfFileExists{xurl.sty}{\usepackage{xurl}}{} % add URL line breaks if available
\urlstyle{same}
\hypersetup{
  pdftitle={COMP3000 Week 1},
  pdfauthor={COMP3000 Programming Languages},
  hidelinks,
  pdfcreator={LaTeX via pandoc}}

\title{COMP3000 Week 1}
\subtitle{Team Based Learning + Domain-Specific Languages}
\author{COMP3000 Programming Languages}
\date{}

\begin{document}
\frame{\titlepage}

\begin{frame}{Welcome to COMP3000}
\phantomsection\label{welcome-to-comp3000}
\begin{itemize}
\tightlist
\item
  13 lectures, 110 minutes each
\item
  Every lecture is followed by a two-hour collaborative workshop
\item
  Today: what the unit is about, how classes run, and the big idea that
  drives everything
\item
  No preparation needed this week -- but from Week 2 on there is
\end{itemize}

\note{Walk through the shape of the unit before diving into content.
Emphasise the lecture-to-workshop rhythm: the lecture introduces, the
workshop applies.}
\end{frame}

\begin{frame}{What the unit is about}
\phantomsection\label{what-the-unit-is-about}
\begin{itemize}
\tightlist
\item
  You will \textbf{design and build your own programming language}
\item
  Not a toy: a language that solves a real problem
\item
  The domain this year: \textbf{modelling water flows in river systems}
\end{itemize}

\begin{center}
\includegraphics[width=0.45\textwidth]{googong.jpg}
\end{center}

\emph{The Googong Dam -- one of the river flows students model.}

\note{Show the dam image. This is the hook: languages are usually
described as general-purpose tools, but here we build a \emph{little
language} for a specific job, the way industry actually uses them.}
\end{frame}

\begin{frame}{Why a language for water flows?}
\phantomsection\label{why-a-language-for-water-flows}
\begin{itemize}
\tightlist
\item
  River managers need to know what rainfall does to river levels over
  time
\item
  Flood prediction, dam strategies, public safety
\item
  Normally done with simulation software
\item
  We will do it with a \textbf{custom language} instead
\item
  No such language exists yet -- you will build it
\end{itemize}

\note{Frame as EPIC: purpose that inspires, impact that matters. The
students are filling a real hole.}
\end{frame}

\begin{frame}{How the classes run -- Team Based Learning}
\phantomsection\label{how-the-classes-run-team-based-learning}
\begin{itemize}
\tightlist
\item
  The unit runs on \textbf{TBL}: you work in a team for all 13 weeks
\item
  Each week follows a fixed cycle:

  \begin{itemize}
  \tightlist
  \item
    \textbf{iRAT} -- individual multiple-choice quiz (done before class)
  \item
    \textbf{tRAT} -- the same quiz, answered with your team
  \item
    \textbf{FAT} -- class-wide discussion of the answers
  \item
    \textbf{Application Task} -- one hour of team work on the assignment
  \end{itemize}
\item
  Pre-reading every week, starting Week 2
\end{itemize}

\note{iRAT before class; tRAT + FAT + application task in class. Team
contracts are made in Week 1.}
\end{frame}

\begin{frame}{Teams matter}
\phantomsection\label{teams-matter}
\begin{itemize}
\tightlist
\item
  You'll form teams today and stay in them all semester
\item
  The workshops build your assignment \emph{as a team}
\item
  Everyone draws on the shared pool of team work for their own
  submission
\item
  Teams write a \textbf{contract} in Week 1: expectations, and how a
  member can be excluded
\item
  If you stop showing up, your team can replace you
\end{itemize}

\note{Stress the team contract activity in today's workshop. The logbook
is personal, but the work is shared.}
\end{frame}

\begin{frame}{Domain-specific languages -- the big idea}
\phantomsection\label{domain-specific-languages-the-big-idea}
\begin{itemize}
\tightlist
\item
  A \textbf{domain-specific language} (DSL) is a language built for one
  job
\item
  Ch 1 calls them \emph{``little languages''} -- ``pidgins tailor-built
  to a specific task''
\item
  Think application scripting languages, template engines, markup
  formats, config files
\item
  Contrast with a \textbf{general-purpose language} (GPL): Java, C,
  Python -- built for any job
\end{itemize}

\note{Definition from the textbook: ``For every successful
general-purpose language, there are a thousand successful niche ones.''}
\end{frame}

\begin{frame}{Little languages are everywhere}
\phantomsection\label{little-languages-are-everywhere}
\includegraphics[width=0.5\textwidth]{text/site/image/introduction/little-languages.png}

\begin{itemize}
\tightlist
\item
  SQL, regex, LaTeX, Make, CSS, shell -- all little languages
\item
  Almost every large software project needs a handful
\item
  When you can, \textbf{reuse an existing one} instead of rolling your
  own
\end{itemize}

\note{``Once you factor in documentation, debuggers, editor support,
syntax highlighting, and all the other trappings, doing it yourself
becomes a tall order.'' -- ch 1}
\end{frame}

\begin{frame}{Some well-known DSLs}
\phantomsection\label{some-well-known-dsls}
\begin{itemize}
\tightlist
\item
  \textbf{SQL} -- querying relational databases
\item
  \textbf{Regular expressions} -- pattern matching in text
\item
  \textbf{LaTeX} -- typesetting documents
\item
  \textbf{Make} -- describing build dependencies
\item
  \textbf{Dream Maker} -- the game \emph{Space Station 13} is built in
  it
\item
  \textbf{strudel.cc} -- making music by writing code in the browser
\end{itemize}

\note{Mention the live regex demo that comes in Week 2. strudel.cc is
worth showing live if internet allows.}
\end{frame}

\begin{frame}{Why build a DSL at all?}
\phantomsection\label{why-build-a-dsl-at-all}
\begin{itemize}
\tightlist
\item
  When a general-purpose language is \textbf{too much}:

  \begin{itemize}
  \tightlist
  \item
    too verbose for the job
  \item
    too hard for the people who need to write it
  \end{itemize}
\item
  A DSL says \emph{exactly} what the task needs, in the task's own terms
\item
  The classic trade-off from ch 1: reuse an existing one\ldots{} or
  build your own when nothing fits
\end{itemize}

\note{Domain experts (here: water engineers) shouldn't need a CS degree
to express their problem.}
\end{frame}

\begin{frame}{Languages are great exercise}
\phantomsection\label{languages-are-great-exercise}
\begin{itemize}
\tightlist
\item
  Implementing a language is a \emph{real} test of programming skill
\item
  Recursion, trees, graphs, dynamic arrays, hash tables -- all from
  scratch
\item
  ``Rise to it, and you'll come away a stronger programmer''
\end{itemize}

\note{From ch 1 ``Languages are great exercise''.}
\end{frame}

\begin{frame}{There's no magic}
\phantomsection\label{theres-no-magic}
\begin{itemize}
\tightlist
\item
  Languages feel magical -- until you build one
\item
  ``There is no magic at all. It's just code, and the people who hack on
  languages are just people.'' -- ch 1
\item
  By the end of this unit you'll \emph{be} one of those people
\end{itemize}

\note{The demystifying pitch. A couple of techniques are new, but
nothing harder than obstacles already overcome.}
\end{frame}

\begin{frame}{The pipeline -- a teaser}
\phantomsection\label{the-pipeline-a-teaser}
\begin{itemize}
\item
  A language implementation transforms source text step by step
\item
  The three stages we build this semester:

  \textbf{Scanner} \(\rightarrow\) \textbf{Parser} \(\rightarrow\)
  \textbf{Interpreter}
\item
  Today, just the names. We build each in turn from Week 4.
\end{itemize}

\note{Ch 2's mountain metaphor: source text at the bottom, climb to
where the code's meaning is clear, descend to something the machine can
run. Our first interpreter is a \emph{tree-walk} interpreter -- climb
partway, then evaluate the tree directly.}
\end{frame}

\begin{frame}[fragile]{Source text becomes tokens}
\phantomsection\label{source-text-becomes-tokens}
\texttt{var\ average\ =\ (min\ +\ max)\ /\ 2;}

\vspace{0.5em}


\includegraphics[width=0.8\textwidth]{text/site/image/a-map-of-the-territory/tokens.png}

\emph{The scanner chunks characters into tokens.}

\note{Scanning (lexing): the linear stream of characters becomes
words/tokens. Whitespace and comments are discarded.}
\end{frame}

\begin{frame}{Tokens become a tree}
\phantomsection\label{tokens-become-a-tree}
\includegraphics[width=0.75\textwidth]{text/site/image/a-map-of-the-territory/ast.png}

\emph{The parser builds an abstract syntax tree from the flat token
stream.}

\note{Parsing gives the syntax a grammar -- the ability to compose
larger expressions from smaller parts, mirroring the nested structure of
the language.}
\end{frame}

\begin{frame}{The mountain, in full}
\phantomsection\label{the-mountain-in-full}
\includegraphics[width=0.9\textwidth]{text/site/image/a-map-of-the-territory/mountain.png}

\begin{itemize}
\tightlist
\item
  Not every implementation walks the whole mountain
\item
  Our interpreter: \textbf{scan, parse, evaluate} -- a tree-walk
  interpreter
\item
  Later topics (Ch 11) add a static-analysis pass before evaluating
\end{itemize}

\note{Ch 2's map: scanning, parsing, static analysis, optimization, code
generation, VM, runtime. We take the tree-walk route this semester.}
\end{frame}

\begin{frame}{Wrap-up and Week 2 prep}
\phantomsection\label{wrap-up-and-week-2-prep}
\begin{itemize}
\tightlist
\item
  Today's reading: \textbf{Chapters 1--2} of \emph{Crafting
  Interpreters}
\item
  \textbf{Before next class}: read those chapters and complete the Week
  2 iRAT
\item
  Next week: little languages in more depth, and the \emph{full}
  pipeline picture
\item
  Reminder: form your team and write your team contract today
\end{itemize}

\note{Next lecture covers the pipeline in full plus regular expressions
as a worked DSL. iRAT is done individually before class, every week from
now on.}
\end{frame}

\end{document}
